% =====================================================================
%  1bat-ud2-3.tex  —  HORA 3: Formalització del llenguatge algebraic
%  (sense preàmbul: s'inclou des de main.tex)
% =====================================================================
\horatitol{Sessió 3 — Formalització del llenguatge algebraic}%
{Vocabulari oficial: terme, coeficient, terme independent, grau i polinomi.}

\subsection{Idea clau}

Una \textbf{expressió algebraica} és una combinació de nombres i lletres unides
per operacions. Quan només hi apareixen sumes, restes i productes de potències
naturals de la variable, parlem d'un \textbf{polinomi}.

\begin{idea}
\textbf{Vocabulari} (sobre l'exemple $\;P = 5x^2 - 3x + 8$):
\begin{itemize}[leftmargin=1.4em, itemsep=2pt]
  \item \textbf{Termes}: $5x^2$, \;$-3x$, \;$8$ (les parts separades per $+$ o $-$).
  \item \textbf{Coeficients}: $5$ (del terme $x^2$) i $-3$ (del terme $x$).
  \item \textbf{Terme independent}: $8$ (no porta lletra).
  \item \textbf{Grau}: l'exponent més gran de la variable; aquí és $2$.
\end{itemize}
\end{idea}

Un polinomi de \textbf{grau 1} (com $C=2000+15x$) modelitza un cost que creix de
manera \emph{constant}. Un polinomi de \textbf{grau 2} (amb un terme $x^2$, com
$C(x)=\num{0,05}x^2-10x+500$) modelitza situacions on el ritme de creixement
canvia, com veurem a la sessió 4.

\subsection{Exemples}

\begin{exemple}[primer grau vs segon grau]
$\;P_1 = 7x + 4\;$ és de \textbf{grau 1} (terme de més grau: $7x$). \\[2pt]
$\;P_2 = 2x^2 + 7x + 4\;$ és de \textbf{grau 2} (terme de més grau: $2x^2$).
\end{exemple}

\begin{exemple}[ordenar un polinomi]
El polinomi $\;3 - 4x + x^2\;$ es reescriu ordenat de més grau a menys:
\[
x^2 - 4x + 3.
\]
Coeficients: $1$ (de $x^2$) i $-4$ (de $x$); terme independent: $3$; grau $2$.
\end{exemple}

\subsection{Exercicis resolts}

\begin{exresolt}[identificar elements]
Donat el polinomi $\;P = -2x^2 + 9x - 5$, indica'n els termes, els coeficients,
el terme independent i el grau.
\solucio
\begin{itemize}[leftmargin=1.4em, itemsep=2pt]
  \item Termes: $-2x^2$, \;$9x$, \;$-5$.
  \item Coeficients: $-2$ (de $x^2$) i $9$ (de $x$). \emph{Atenció al signe.}
  \item Terme independent: $-5$.
  \item Grau: $2$ (l'exponent més gran).
\end{itemize}
\resposta{Grau $2$; coeficients $-2$ i $9$; terme independent $-5$.}
\end{exresolt}

\begin{exresolt}[reescriure una regla amb notació formal]
A la sessió 2 vam trobar, per a un centre, la regla «cost fix $\num{2000}$ més
$15$ euros per usuari, i a més un terme de $\num{0,5}$ euros per usuari al quadrat
en concepte de logística». Escriu el polinomi corresponent, ordenat, i indica'n el grau.
\solucio
La part «per usuari al quadrat» és $\num{0,5}x^2$; la «per usuari» és $15x$;
la fixa és $\num{2000}$. Ordenat de més grau a menys:
\[
C(x) = \num{0,5}\,x^2 + 15x + \num{2000}.
\]
És un polinomi de \textbf{grau 2}, perquè el terme de més grau és $\num{0,5}x^2$.
\resposta{$C(x)=\num{0,5}x^2+15x+\num{2000}$, de grau $2$.}
\end{exresolt}

\subsection{Exercicis per fer}

\begin{exercicis}
\begin{exlist}
  \item Per a cada polinomi, indica els coeficients, el terme independent i el grau:
    \begin{enumerate}[label=\alph*), leftmargin=1.6em, topsep=2pt]
      \item $P = 4x + 9$
      \item $Q = x^2 - 6x + 2$
      \item $R = -3x^2 + 5$
      \item $S = 10x$
    \end{enumerate}
  \item Ordena de més grau a menys i indica'n el grau:
    \quad (a) $7 - 2x + x^2$ \quad (b) $-x + 4 + 3x^2$.
  \item Classifica cada expressió com a polinomi de \emph{grau 1} o de \emph{grau 2}:
    \quad (a) $8x-1$ \quad (b) $x^2+x$ \quad (c) $\num{0,2}x^2-3x+7$ \quad (d) $25x+100$.
  \item Escriu un polinomi de grau $1$ amb coeficient $6$ i terme independent $-4$.
  \item Escriu un polinomi de grau $2$ amb coeficient de $x^2$ igual a $\num{0,1}$,
        coeficient de $x$ igual a $-8$ i terme independent $400$.
  \item \textbf{(Context)} Una entitat modelitza la despesa d'un projecte amb
        $D(x) = \num{0,3}x^2 + 20x + \num{1500}$, on $x$ és el nombre de
        beneficiaris. Indica el grau del polinomi i explica què representa,
        en aquest cas, el terme independent $\num{1500}$.
\end{exlist}
\end{exercicis}

% ---------------------------------------------------------------------
\fullsolucions{Sessió 3}
\begin{solucions}
\begin{sollist}
  \item (a) coeficient $4$, terme independent $9$, grau $1$ \quad
        (b) coeficients $1$ i $-6$, terme independent $2$, grau $2$ \quad
        (c) coeficient $-3$ (de $x^2$), terme independent $5$, grau $2$ \quad
        (d) coeficient $10$, terme independent $0$, grau $1$
  \item (a) $x^2 - 2x + 7$, grau $2$ \quad (b) $3x^2 - x + 4$, grau $2$
  \item (a) grau $1$ \quad (b) grau $2$ \quad (c) grau $2$ \quad (d) grau $1$
  \item $6x - 4$
  \item $\num{0,1}\,x^2 - 8x + 400$
  \item Grau $2$. El terme independent $\num{1500}$ és la despesa fixa del
        projecte (la que hi ha encara que no s'atengui cap beneficiari, $x=0$).
\end{sollist}
\end{solucions}
